% !TeX program = xelatex
% Progressive Frequency-conditioned FiLM for iTransformer-based
% Long-term Time Series Forecasting
% Compile (recommended):
%   latexmk -xelatex -interaction=nonstopmode main.tex

\documentclass[UTF8,a4paper,11pt]{ctexart}

% -------------------------
% Packages
% -------------------------
\usepackage{amsmath,amssymb,bm}
\usepackage{booktabs,multirow}
\usepackage{graphicx}
\usepackage{subcaption}
\usepackage{siunitx}
\usepackage{geometry}
\usepackage{microtype}
\usepackage{enumitem}
\usepackage{hyperref}

\geometry{margin=1in}
\setlist[itemize]{noitemsep,topsep=2pt,leftmargin=1.25em}
\setlist[enumerate]{noitemsep,topsep=2pt,leftmargin=1.25em}
\hypersetup{
  colorlinks=true,
  linkcolor=blue,
  citecolor=blue,
  urlcolor=blue
}

% -------------------------
% Title & Author (placeholders)
% -------------------------
\title{Progressive Frequency-conditioned FiLM for Transformer-based Long-term Time Series Forecasting}
\author{作者A\thanks{单位与邮箱占位} \and 作者B\thanks{单位与邮箱占位}}
\date{} % leave empty for journal-style

\begin{document}
\maketitle

% ============================================================
% 旧版摘要（已注释）
% ============================================================
% 多变量时间序列预测与基于观测序列的物理量回归在工程与科学计算中均十分关键：前者服务于能源、交通与工业系统的未来态估计，后者可用于由观测序列反演天体/材料等对象的关键参数（例如 CEMP 恒星参数回归）。两类任务共享跨维耦合、长程依赖与显著非平稳性等挑战，但输出形式可分别对应多步预测（sequence-to-sequence）与多变量回归（sequence-to-vector）。
%
% 本文提出渐进式频率条件 FiLM（Progressive Frequency-conditioned FiLM），一种轻量且高效的统一序列建模框架，以 iTransformer 为骨干，通过逐层渐进衰减的特征线性调制将多分辨率频谱信息注入 Transformer Encoder。其核心设计包含三个要素：...
%
% 在公开预测数据集与 CEMP 参数回归数据上的实验（待补充）表明，该方法在多种设置下取得一致提升；消融实验进一步验证了多分辨率频谱、渐进式调制与可学习强度的各自贡献。

% ============================================================
% 新版摘要（按顶刊结构重写）
% ============================================================
\begin{abstract}
时间序列预测中利用频域信息的现有方法大致可分为两类：频域原生模型将注意力或特征聚合操作迁移至频域执行，具有较强的频谱建模能力，但通常需要对网络结构进行较大改动，难以作为轻量模块嵌入现有骨干模型；频域条件模型则以附加特征形式注入频谱信息，具有良好的架构兼容性，但大多仅采用单点、固定强度的调制方式，深层传播过程中频域引导易被逐渐稀释，且调制强度依赖人工设定，缺乏跨数据集的自适应能力。

为弥合上述两类方法之间的鸿沟，本文提出一种轻量级频域条件调制框架——\textbf{渐进式频率条件 FiLM}（Progressive Frequency-conditioned FiLM）。该方法无需修改主干网络结构，而是通过数据驱动方式感知序列的频谱状态，并以自适应强度逐层调控 Transformer Encoder 的隐藏表征，使频域信息能够贯穿编码全过程，同时在缺乏显著频谱结构的数据上自动抑制调制以避免有害注入。

具体而言，本文从三个方面提升频域条件建模能力：其一，设计多分辨率频谱感知机制，从不同时间尺度刻画序列的周期结构与可预测性；其二，引入逐层渐进式频域调制策略，使频域引导能够在编码过程中持续传播，并随网络深度动态衰减，从而兼顾全局谱结构建模与深层语义表达；其三，提出可学习的调制强度控制机制，使模型能够依据数据特性自适应决定频域信息的介入程度，而非依赖人工设定的固定调制系数。理论分析与实验结果表明，该方法能够在不破坏主干网络收敛性的前提下扩展模型的可表达空间，并以极低额外计算开销，在多变量长时预测与 CEMP 恒星参数回归等任务上取得一致且显著的性能提升。
\end{abstract}

\noindent\textbf{关键词：} 长时序列预测；频率条件 FiLM；渐进式调制；多分辨率频谱；非平稳性；CEMP

% 01 Introduction
\section{引言}
\label{sec:introduction}

% 序列数据广泛存在于工程系统与科学观测中，是刻画复杂动态过程与观测机制的核心载体。时间序列记录了系统随时间变化的观测值，因而在制造、物流、能源与医疗保健等多个领域至关重要：例如预测电力消耗、交通流量、产品销量与库存水平，或对工业过程进行早期预警与资源调度。在科学观测中，序列数据同样扮演关键角色：诸如由光谱/测光等观测序列反演天体的物理与化学属性，可为恒星演化与银河化学演化研究提供基础支撑，其中 CEMP（Carbon-Enhanced Metal-Poor）恒星的参数回归是具有代表性的应用场景。从学习形式上看，这类问题可统一表述为\emph{序列到多变量输出}学习：既包括多变量时间序列预测（MTSF）的多步预测（sequence-to-sequence），也包括参数反演的多变量回归（sequence-to-vector）。尽管输出形态不同，两者都需要在噪声与观测不完备条件下刻画跨维耦合与长程依赖，并在分布漂移与统计量变化下保持稳定泛化。

序列数据广泛存在于工程系统与科学观测中，是刻画复杂动态过程与观测机制的核心载体。时间序列记录了系统随时间变化的观测值，因而在制造、物流、能源与医疗保健等多个领域至关重要：例如预测电力消耗、交通流量、产品销量与库存水平，或对工业过程进行早期预警与资源调度。在科学观测中，序列数据同样扮演关键角色：诸如由光谱/测光等观测序列反演天体的物理与化学属性，可为恒星演化与银河化学演化研究提供基础支撑，其中 CEMP（Carbon-Enhanced Metal-Poor）恒星的参数回归是具有代表性的应用场景。尽管上述任务在目标形式上分别对应未来序列生成与物理参数估计，但其本质都属于基于序列观测的多变量建模问题，即模型不仅需要从复杂时序结构中捕获跨变量耦合关系与长程依赖，还需在噪声干扰、观测不完备以及数据分布变化等条件下保持稳定的泛化能力。

基于从传统统计模型到深度模型（如 RNN/TCN），特别是近年来基于 MLP 和 Transformer 架构的研究进展，加上输入表示层面（如 Patch 分块与协变量引入）的持续优化，已经迅速提高了时间序列预测的精度。然而，现有的深度网络在处理复杂多变量任务时，往往依赖隐式的频率状态利用和单一的结构先验，使其在非平稳条件下的动态适应能力严重受限。这种缺乏协同动态调制与稳定化机制的缺陷在实际应用中是具有不确定性的。在具有高噪声和异质数据的现实场景中，系统不太可能依赖那些难以应对分布漂移且极易出现训练不稳的预测模型。在对泛化稳定性和可靠性至关重要的领域部署此类鲁棒性不足的模型会带来重大风险。如果能在高精度的同时显式引入与时序信号一致的频率归纳偏置并协同跨变量建模，模型就能够有效应对极端变化。这一主题被称为非平稳时序的鲁棒建模，并在复杂场景的时间序列预测与参数回归领域被积极研究。一些现有方法试图通过降低长序列注意力复杂度（Informer~\cite{informer}）、从分解与频域角度增强周期结构刻画（Autoformer~\cite{autoformer}; FEDformer~\cite{fedformer}）、通过变量维 token 化提升跨变量交互能力（iTransformer~\cite{itransformer}）以及针对分布漂移提供可逆归一化机制（RevIN~\cite{revin}）来应对这些挑战。然而，如何在高精度的同时引入与时序信号一致的频率归纳偏置，并在非平稳条件下保持训练与泛化稳定性，仍是模型走向实际应用的关键问题。

为解决上述问题，本文提出\textbf{渐进式频率条件 FiLM}（Progressive Frequency-conditioned FiLM），一种以 Transformer 为骨干、通过逐层渐进衰减的 FiLM 将多分辨率频谱信息注入 Transformer Encoder 的轻量高效框架。其核心创新包含三个要素：

\textbf{(1) 多分辨率频谱提取（Multi-Resolution Spectrum）}：在三个时间窗口（全序列 \(L\)、半序列 \(L/2\)、四分之一序列 \(L/4\)）上分别通过 FFT 提取功率谱，每窗口独立计算多频带归一化能量、高低频比值与频谱熵，形成尺度鲁棒且信息丰富的频域表征。相比固定低通滤波，FFT 频谱能够更完整地捕获周期结构与谱平坦度等可预测性信号。

\textbf{(2) 逐层渐进式 FiLM（Layer-wise Progressive FiLM）}：压缩后的频谱 latent 通过每层独立的小 MP 生成调制参数 \((\alpha, \beta)\)，注入到每一个 Transformer Encoder Layer 之前；调制强度按 \(\gamma_{\text{base}} \cdot \delta^{\ell}\) 随层深指数衰减（\(\delta \in (0,1)\) 为衰减因子），实现由粗到细的层级化频率条件调控——浅层接受较强的频域引导以建立全局谱结构感知，深层逐渐减弱调制以保护语义表征的精细辨别能力。

\textbf{(3) 可学习调制强度（Learnable Modulation Scale）}：将全局调制基强度 \(\gamma_{\text{base}}\) 参数化为 \(\mathrm{sigmoid}(\theta_{\gamma})\)（其中 \(\theta_{\gamma}\) 为可学习参数，初始对应 \(\gamma_{\text{base}} \approx 0.02\)），替代手工设定的固定系数。不同数据集在训练中自适应决定频域信息的注入程度，同时兼顾可退化性：初始状态下调制接近关闭，模型等价于纯 iTransformer。

上述设计与 iTransformer~\cite{itransformer} 的变量维 token 化策略及 RevIN~\cite{revin} 可逆归一化无缝协同，并通过 FiLM 末层零初始化、\(\tanh\) 有界约束与 \(\log(1+x)\) 压缩等稳定化设计保证训练稳健。频率条件信号具有明确物理/信号含义（频带能量分布、高低频能量对比、频谱平坦度），为理解模型在不同序列形态下的响应差异提供了可分析性。框架既可用于 MTSF 多步预测，也可通过更换任务头用于 CEMP 等场景下的多变量参数回归。

\subsection{本文贡献}
\label{subsec:contributions}
本文的主要贡献概括如下（可根据最终实验进一步精炼/增补）：
\begin{itemize}
  \item 提出\textbf{多分辨率频谱特征提取}方法，以 FFT 功率谱的多频带归一化能量、高低频比值与频谱熵联合刻画输入谱状态，相比传统低通滤波方案具有更强的尺度鲁棒性与信息完整性。
  \item 设计\textbf{逐层渐进式 FiLM 调制}机制：每层独立生成调制参数并以指数衰减强度逐层注入，在不破坏 iTransformer 主干可训练性的前提下实现层级化频域引导。
  \item 提出\textbf{可学习调制强度}参数化方案（Sigmoid 门控 + 保守初始化），使模型自适应决定频域信息的注入程度，在无频域信号的数据上可自动退化为纯 iTransformer。
  \item 给出\textbf{面向非平稳与训练稳定的整体方案}：RevIN 处理分布漂移，配合频谱压缩、\(\log(1+x)\) 压缩/截断、\(\tanh\) 有界调制与零初始化，增强鲁棒性与可训练性。
  \item 在\textbf{多变量预测与 CEMP 参数回归}两类任务上的实验与消融分析（待补充）验证方法有效性，并量化各模块对性能的贡献。
\end{itemize}

\subsection{论文结构}
本文余下部分组织如下：第\ref{sec:related_work}节介绍背景与相关工作；
第\ref{sec:method}节给出方法细节；
第\ref{sec:experiments}节描述实验设置并报告结果；
第\ref{sec:conclusion}节总结全文并讨论未来工作。

\section{相关工作}
\label{sec:related_work}

\subsection{时间序列预测的深度学习方法}
时间序列研究大体经历了从统计方法（ARIMA~\cite{arima_survey}、VAR~\cite{var_survey}）到深度方法、再到 Transformer 体系的演进。围绕长序列建模效率，Informer 通过 ProbSparse 注意力降低计算复杂度~\cite{informer}；围绕结构先验注入，Autoformer 利用趋势--季节分解与自相关机制建模周期结构~\cite{autoformer}，FEDformer 将频域表示引入全局模式学习~\cite{fedformer}；iTransformer 则创新性地将传统"时间步 token 化"反转为"变量维 token 化"，以 Transformer Encoder 建模跨变量交互而非跨时间步依赖，在多变量预测任务上取得了显著提升~\cite{itransformer}。与此同时，patch 分块处理、协变量融合与线性模型等方法也持续推动性能边界。本文以 iTransformer 为骨干架构，保留其变量维 token 化与 Encoder 设计，并通过频谱条件调制增强其频率感知能力。

\subsection{频域分析与多分辨率表征}
频域方法在时间序列建模中具有重要的归纳偏置价值。FEDformer 通过傅里叶/小波变换在频域执行注意力操作~\cite{fedformer}；TimesNet 将一维序列变换为二维张量以捕获多周期模式；相关频谱分析方法利用 FFT 提取功率谱、相位谱等特征进行序列分类与异常检测~\cite{fourier_ts}。然而，多数现有方法或完全在频域内操作（引入较大架构改动），或仅将频谱作为辅助可视化工具，缺乏将其作为轻量条件信号注入通用骨干网络的机制。

多分辨率/多尺度建模通过在不同时间粒度上聚合信息以增强表征鲁棒性：金字塔结构融合多层级特征~\cite{ts_pyramid}，多尺度卷积/池化并行的思路在计算机视觉与序列任务中均有应用~\cite{multiscale_ts}。但在时序 Transformer 中，多尺度设计常伴随计算开销上升与架构复杂化等问题。本文的多分辨率频谱在概念上不同于时域多尺度池化——我们通过在不同长度的时间窗口上独立计算 FFT，获得"序列长度鲁棒"的频谱特征，尺度数量固定（3 个）且计算量可控；频带能量归一化进一步保证了不同窗口特征的可比性。

\subsection{条件调制与特征归一化}
条件调制（Conditional Modulation）通过学习输入依赖的缩放/平移参数实现样本自适应表征调整。FiLM（Feature-wise Linear Modulation）最早在视觉推理中提出~\cite{film}，后续在风格迁移、少样本学习与序列建模中得到了广泛应用~\cite{cond_norm,dyn_conv,adaptive_norm_ts}。在时间序列领域，部分工作开始探索以频率特征作为调制条件。然而，直接将 FiLM 应用于深度 Transformer 存在两个风险：(1) 调制强度过大可能淹没原始表征，破坏骨干网络的可训练性；(2) 单一调制点（仅作用于输入端）在深层网络中梯度衰减严重，调制效果逐层稀释。

本文的逐层渐进式 FiLM 直接针对上述问题：(1) 在每个 Encoder Layer 前独立注入调制，避免深层梯度衰减；(2) 调制强度按 \(\gamma_{\text{base}} \cdot \delta^{\ell}\) 指数衰减，浅层接受较强频域引导以建立全局结构感知，深层逐渐减弱以保护语义表征；(3) 可学习 gamma 与零初始化保证初始等价于无调制的骨干网络，频域信息在训练中逐步激活。该设计受稳定门控（stable gating）相关讨论的启发~\cite{stable_gating}，但将其系统化为一种有理论直觉（由粗到细）的层级化调制范式。

\subsection{非平稳性与分布漂移处理}
非平稳性是现实世界时间序列的核心挑战之一。RevIN（Reversible Instance Normalization）通过在归一化--反归一化框架下对齐每个样本的统计量，以极低的计算开销有效缓解了分布漂移对预测模型的影响~\cite{revin}，已成为许多现代时序模型的标准组件。然而 RevIN 本身不建模信号的频率特征，本文在保留 RevIN 的同时，通过 FFT 频谱条件为模型补充了归一化过程中可能被削弱的频域信息。

\subsection{科学观测序列的参数回归}
以 CEMP 恒星参数回归为代表的任务旨在从光谱/测光序列中反演 \(T_{\mathrm{eff}}\)、\(\log g\)、\([\mathrm{Fe/H}]\)、\([\mathrm{C/Fe}]\) 等物理参数。相关工作中，StarNet 采用卷积神经网络提取局部谱特征~\cite{starnet}，但在受红化、仪器差异或低信噪比影响的非平稳观测条件下，由于缺乏全局频率感知的动态调制机制，泛化鲁棒性仍有提升空间。本文统一框架通过更换任务头（预测头 \(\to\) 回归头）即可适配此类 sequence-to-vector 任务，频谱条件信号为模型提供了信噪比自适应的感知能力。

综合来看，现有工作对"多变量预测与科学观测参数回归"的统一场景，仍缺少一种同时兼顾 FFT 频谱编码、层级化自适应调制与可退化训练的轻量化框架，本文方法即在该研究空白上展开。

% 03 Method
\section{方法}
\label{sec:method}

本节详细描述渐进式频率条件 FiLM 框架的整体架构与关键模块。为便于表述，设输入序列为 \(X \in \mathbb{R}^{B \times L \times C}\)，其中 \(B\) 为 batch 大小，\(L\) 为历史序列长度（回望窗口），\(C\) 为变量维度（通道数）。模型整体遵循 iTransformer~\cite{itransformer} 的变量维 token 化范式，将每个变量的长度为 \(L\) 的历史片段映射为一个 \(d_{\text{model}}\) 维 token，并通过 Transformer Encoder 建模跨变量交互。本文的核心贡献在于：(1) 在变量 token 流入 Encoder 之前及每一层之前，注入由 FFT 多分辨率频谱生成的条件调制信号；(2) 调制采用逐层独立生成、强度随深度指数衰减的渐进式 FiLM 机制；(3) 全局调制基强度为可学习参数，支持数据集自适应与可退化训练。

\subsection{总体架构}
\label{subsec:overview}

渐进式频率条件 FiLM 的前向计算流程如下（见图\ref{fig:architecture}，待补充）：

\begin{enumerate}[leftmargin=*,label=\textbf{Step \arabic*:},itemsep=4pt]
  \item \textbf{RevIN 归一化}：对输入 \(X\) 沿时间维做可逆实例归一化，得到 \(\tilde{X} \in \mathbb{R}^{B \times L \times C}\)，缓解非平稳分布漂移；
  \item \textbf{变量 Token 嵌入}：将每个变量的 \(L\) 步序列线性映射为 \(d_{\text{model}}\) 维 token，并通过 LayerNorm，得到 \(T^{(0)} \in \mathbb{R}^{B \times C \times d_{\text{model}}}\)；
  \item \textbf{多分辨率频谱提取}：在三个时间窗口上分别对 \(\tilde{X}\) 做 FFT 并提取频带能量、高低频比与频谱熵，得到多分辨率频谱特征 \(S \in \mathbb{R}^{B \times C \times D_s}\)；
  \item \textbf{频谱压缩}：通过小型 MLP 将高维频谱特征压缩为低维 latent \(Z \in \mathbb{R}^{B \times C \times d_z}\)，防止频域信息过载；
  \item \textbf{逐层渐进式 FiLM 调制 + Transformer 编码}：对 \(\ell = 0, 1, \dots, N-1\) 层，依次执行：
  \begin{enumerate}
    \item 由第 \(\ell\) 层的独立 FiLM Generator 从 \(Z\) 生成调制参数 \((\alpha_\ell, \beta_\ell)\)；
    \item 以强度 \(\gamma_\ell = \gamma_{\text{base}} \cdot \delta^{\ell}\) 对 token 做有界 FiLM 调制；
    \item 经 Pre-LN 与 Transformer Encoder Layer 更新 token 表征，得到 \(T^{(\ell+1)}\)。
  \end{enumerate}
  \item \textbf{预测头与 RevIN 反归一化}：将最终 token 表征经 LayerNorm 与线性映射得到未来 \(H\) 步预测，再通过 RevIN 反归一化恢复到原始尺度。
\end{enumerate}

对于 CEMP 参数回归等 sequence-to-vector 任务，Step 6 中的预测头替换为回归头（token 维池化后接 MLP 映射到参数向量），且根据任务需要可选择是否执行 RevIN 反归一化。

\subsection{RevIN：可逆实例归一化}
\label{subsec:revin}

RevIN~\cite{revin} 对每个样本沿时间维独立计算均值与标准差，将序列归一化到零均值单位方差的分布，并在输出端通过反归一化恢复原始尺度。其形式为：

\begin{align}
  \tilde{X} &= \gamma \odot \frac{X - \mu}{\sigma + \epsilon} + \beta, \quad
  \mu = \frac{1}{L}\sum_{t=1}^{L} X_{:,t,:},\;
  \sigma = \sqrt{\frac{1}{L}\sum_{t=1}^{L} (X_{:,t,:} - \mu)^2 + \epsilon}, \\
  \hat{Y} &= \frac{\hat{Y}_{\text{raw}} - \beta}{\gamma} \odot \sigma + \mu,
\end{align}

其中 \(\gamma, \beta \in \mathbb{R}^{C}\) 为可学习的仿射参数，\(\mu, \sigma\) 在归一化时 detach（不参与梯度回传）。通过 RevIN，模型的骨干部分始终在统计量对齐的分布上学习，有效降低了非平稳性对训练与泛化的干扰。

\subsection{变量 Token 嵌入}
\label{subsec:token_embed}

遵循 iTransformer 的设计，将每个变量的整个历史序列片段作为一个 token：
\begin{equation}
  T^{(0)} = \mathrm{LayerNorm}\big(\phi(\tilde{X}^{\mathsf{T}})\big),
  \quad T^{(0)} \in \mathbb{R}^{B \times C \times d_{\text{model}}},
\end{equation}
其中 \(\phi: \mathbb{R}^{L} \to \mathbb{R}^{d_{\text{model}}}\) 为单层线性映射（无激活函数），\(\tilde{X}^{\mathsf{T}} \in \mathbb{R}^{B \times C \times L}\) 为转置后的归一化序列。该设计将计算复杂度从时间维注意力的 \(O(L^2)\) 降为变量维注意力的 \(O(C^2)\)，尤其适合变量数适中的多变量预测场景。

\subsection{多分辨率频谱特征提取}
\label{subsec:multires_spectrum}

本文提出多分辨率频谱提取器（Multi-Resolution Spectrum Extractor），在三个时间尺度上独立计算 FFT 功率谱并提取结构化的频域特征。该模块的输入为归一化后的序列 \(\tilde{X}^{\mathsf{T}} \in \mathbb{R}^{B \times C \times L}\)，输出为多分辨率频谱特征 \(S \in \mathbb{R}^{B \times C \times D_s}\)。

\medskip\noindent\textbf{多时间窗口.}
在三个递减的时间窗口上分别截取输入序列的末尾部分（最近的信息最重要）：
\begin{equation}
  \mathcal{W} = \{L,\; L/2,\; L/4\},
\end{equation}
对于每个窗口大小 \(w \in \mathcal{W}\)（实际实现中取 \(w \gets \max(8, \lfloor w \rfloor)\)），取 \(X_w = \tilde{X}_{:, :, -w:}^{\mathsf{T}} \in \mathbb{R}^{B \times C \times w}\)。

\medskip\noindent\textbf{功率谱计算.}
对每个窗口的每个变量独立计算实值 FFT（rFFT），并取幅值平方得到功率谱：
\begin{equation}
  \mathcal{F}_w = \mathrm{rFFT}(X_w), \quad
  M_w = |\mathcal{F}_w|^2 \;\in\; \mathbb{R}^{B \times C \times f_w},
\end{equation}
其中频率 bin 数 \(f_w = \lfloor w/2 \rfloor + 1\)。

\medskip\noindent\textbf{每窗口的结构化频域特征.}
对每个窗口的功率谱 \(M_w\)，提取以下三类特征：

\emph{(a) 多频带归一化能量.} 将频率轴均分为 \(K\) 个频带（本文取 \(K = 4\)），计算每频带的平均功率并做 L1 归一化：
\begin{equation}
  E_k^{(w)} = \frac{1}{|\mathcal{B}_k|} \sum_{f \in \mathcal{B}_k} M_w[:,:,f], \quad
  \hat{E}_k^{(w)} = \frac{E_k^{(w)}}{\sum_{j=1}^{K} E_j^{(w)} + \epsilon}, \quad k=1,\dots,K.
\end{equation}
归一化频带能量刻画了信号的功率在不同频率区间的分布形态（例如低频主导 vs 高频主导），且尺度不变。

\emph{(b) 高低频比值.} 取最高频带与最低频带的能量比：
\begin{equation}
  R^{(w)} = \frac{E_K^{(w)}}{E_1^{(w)} + \epsilon}.
\end{equation}
该比值直接反映信号的高频分量相对于低频分量的强弱，是刻画信号"波动性 vs 趋势性"的紧凑归纳偏置。

\emph{(c) 归一化频谱熵.} 将功率谱归一化为概率分布后计算信息熵：
\begin{equation}
  p_f^{(w)} = \frac{M_w[:,:,f]}{\sum_{f'} M_w[:,:,f'] + \epsilon}, \quad
  H^{(w)} = -\frac{\sum_f p_f^{(w)} \log(p_f^{(w)} + \epsilon)}{\log f_w + \epsilon}.
\end{equation}
频谱熵衡量功率在频率轴上的集中程度：低熵对应强周期性/窄带信号（可预测性高），高熵对应白噪声/平坦谱（可预测性低）。除以 \(\log f_w\) 进行归一化使得不同窗口大小的熵值可比。

\medskip\noindent\textbf{特征拼接与数值稳定化.}
将三个窗口的 \(K+2\) 维特征沿最后一维拼接，得到 \(D_s = 3 \times (K+2)\) 维的原始频谱特征：
\begin{equation}
  S_{\text{raw}} = [\hat{E}^{(L)}, R^{(L)}, H^{(L)},\; \hat{E}^{(L/2)}, R^{(L/2)}, H^{(L/2)},\; \hat{E}^{(L/4)}, R^{(L/4)}, H^{(L/4)}].
\end{equation}
随后应用 \(\log(1+x)\) 压缩与截断以保证数值稳定性：
\begin{equation}
  S = \mathrm{clamp}\big(\log(1 + S_{\text{raw}}),\; 0,\; s_{\max}\big), \quad s_{\max}=5.0.
\end{equation}

\subsection{频谱压缩器}
\label{subsec:spectrum_compressor}

为防止高维频谱特征直接注入 Encoder 造成信息过载，我们在频谱特征与 FiLM 生成器之间插入一个频谱压缩器（Spectrum Compressor），将 \(D_s\) 维特征压缩到远小于 \(d_{\text{model}}\) 的低维 latent 空间：
\begin{equation}
  Z = \mathrm{Compressor}(S), \quad
  Z \in \mathbb{R}^{B \times C \times d_z},
\end{equation}
其中 \(\mathrm{Compressor}(\cdot)\) 为一个两层 MLP（带 GELU 激活与 Dropout）：
\begin{equation}
  \mathrm{Compressor}(S) = W_2 \cdot \mathrm{Dropout}\big(\mathrm{GELU}(W_1 \cdot S + b_1)\big) + b_2,
\end{equation}
隐层维度设为 \(\max(64, D_s)\)，latent 维度 \(d_z\) 默认取 32（可配置），显著小于 \(d_{\text{model}}\)（通常为 128--512）。

\subsection{逐层渐进式 FiLM 调制}
\label{subsec:progressive_film}

这是本文最核心的设计——将频域 latent \(Z\) 逐层独立地注入到 Transformer 的每一层之前，并以衰减的强度实现由粗到细的层级化调控。

\medskip\noindent\textbf{可学习调制基强度.}
全局调制强度不由手工设定，而是通过一个可学习的无约束参数 \(\theta_{\gamma}\) 经 Sigmoid 映射得到：
\begin{equation}
  \gamma_{\text{base}} = \sigma(\theta_{\gamma}) = \frac{1}{1 + e^{-\theta_{\gamma}}},
\end{equation}
其中 \(\theta_{\gamma}\) 初始化为 \(\theta_{\gamma}^{(0)} = \ln\frac{0.02}{1-0.02} \approx -3.89\)，使得 \(\gamma_{\text{base}} \approx 0.02\)（保守初始化）。训练过程中，不同数据集自适应学习 \(\theta_{\gamma}\)——若频域信息有效，\(\gamma_{\text{base}}\) 增大；若无效，\(\gamma_{\text{base}}\) 保持较小甚至进一步减小，模型退化为接近纯 iTransformer。

\medskip\noindent\textbf{逐层调制强度衰减.}
对第 \(\ell\) 层（\(\ell = 0, 1, \dots, N-1\)），该层实际使用的调制强度为：
\begin{equation}
  \gamma_{\ell} = \gamma_{\text{base}} \cdot \delta^{\ell},
\end{equation}
其中 \(\delta \in (0, 1)\) 为衰减因子（默认 0.7）。这一设计体现了"由粗到细"的直觉：浅层接受更强的频域引导以建立全局谱结构感知，深层逐渐减弱调制以保护已形成的语义表征不被过度扰动。

\medskip\noindent\textbf{逐层独立 FiLM Generator.}
每一层拥有独立的 FiLM Generator——一个小型 MLP，将共享的频域 latent \(Z\) 映射为该层专属的调制参数：
\begin{equation}
  (\alpha_{\ell}, \beta_{\ell}) = \mathrm{FiLMGen}_{\ell}(Z), \quad
  \alpha_{\ell}, \beta_{\ell} \in \mathbb{R}^{B \times C \times 1}.
\end{equation}

每个 \(\mathrm{FiLMGen}_{\ell}\) 的结构为三层 MLP（含 GELU 激活与 Dropout）：
\begin{equation}
  \mathrm{FiLMGen}(Z) = W_3 \cdot \mathrm{GELU}\big(W_2 \cdot \mathrm{GELU}(W_1 \cdot Z + b_1) + b_2\big) + b_3,
\end{equation}
输出维度为 2（分别对应 \(\alpha\) 和 \(\beta\)），关键设计是最后一层 \(W_3, b_3\) 初始化为零，使得初始状态下 \(\alpha_{\ell} = \beta_{\ell} = 0\)，调制对 token 无影响。

\medskip\noindent\textbf{有界 FiLM 调制.}
在第 \(\ell\) 层 Encoder 之前，对 token 施加经过幅度约束的 FiLM 调制：
\begin{equation}
  T_{\text{mod}}^{(\ell)} = T^{(\ell)} \odot \big(1 + \gamma_{\ell} \cdot \tanh(\alpha_{\ell})\big) + \gamma_{\ell} \cdot \beta_{\ell}.
\end{equation}
其中 \(\tanh(\cdot)\) 将 \(\alpha\) 的缩放效应限制在 \((-1, 1)\) 区间内，乘以小系数 \(\gamma_{\ell}\) 后实际调制幅度在 \((-\gamma_{\ell}, \gamma_{\ell})\) 之间——例如初始状态下约 \((-0.02, 0.02)\)，保证调制仅做"微调"。

\medskip\noindent\textbf{完整逐层计算.}
\begin{equation}
  T^{(\ell+1)} = \mathrm{EncoderLayer}_{\ell}\Big(
    \mathrm{LayerNorm}_{\ell}\big(T_{\text{mod}}^{(\ell)}\big)
  \Big), \quad \ell = 0, 1, \dots, N-1.
\end{equation}
每个 \(\mathrm{EncoderLayer}_{\ell}\) 为标准 Transformer Encoder Layer（Pre-LN 模式，batch\_first=True），包含多头自注意力（变量间交互）与前馈网络。

\subsection{预测头与任务适配}
\label{subsec:head}

经过 \(N\) 层编码后，最终 token 表征经 LayerNorm 与线性映射输出预测：
\begin{equation}
  \hat{Y}_{\text{raw}} = \mathrm{Head}\big(\mathrm{LayerNorm}(T^{(N)})\big),
  \quad \mathrm{Head}: \mathbb{R}^{d_{\text{model}}} \to \mathbb{R}^{H},
\end{equation}
经维度变换后得到 \(\hat{Y}_{\text{raw}} \in \mathbb{R}^{B \times H \times C}\)，再由 RevIN 反归一化恢复到原始尺度。

对于 CEMP 参数回归任务，Head 替换为：Token 维平均池化 \(\to\) 两层 MLP（含 GELU）\(\to\) 参数向量 \(\hat{\bm{p}} \in \mathbb{R}^{P}\)，不执行 RevIN 反归一化。

\subsection{可退化性与稳定化设计总结}
\label{subsec:stability}

渐进式频率条件 FiLM 框架的可退化性由以下三项设计联合保证，三者共同确保模型在初始状态等价于纯 iTransformer，频域调制在训练中"逐步激活"而非"骤然介入"：

\begin{enumerate}[leftmargin=*,itemsep=2pt]
  \item \textbf{FiLM 末层零初始化}：所有 \(\mathrm{FiLMGen}_{\ell}\) 的最后一层 \(W_3, b_3 = 0\)，初始 \(\alpha = \beta = 0\)，调制为恒等操作；
  \item \textbf{保守的 gamma 初始化}：\(\gamma_{\text{base}}\) 初始约 0.02（比常见硬编码值 0.1 更保守），乘以 \(\tanh\) 约束后实际调制幅度极小；
  \item \textbf{数值稳定化管道}：\(\log(1+x)\) 压缩消除长尾分布极端值，clamp 截断防止数值溢出，频谱特征始终在一个可控范围内。
\end{enumerate}

当 \(\gamma_{\text{base}} \to 0\) 时（模型认为频域信息无用），所有层的调制趋近于零，框架严格退化为 iTransformer + RevIN。

% 04 Experiments (placeholders; fill in with actual results)
\section{实验}
\label{sec:experiments}

本节给出实验设置与结果。你可以在补充数据与训练细节后，直接替换占位内容并保留排版结构。

\subsection{数据集}
本工作覆盖两类数据：
\begin{itemize}
  \item \textbf{多变量时间序列预测数据}：填写数据集名称、变量数、采样频率、训练/验证/测试划分方式、缺失值处理与标准化策略等；
  \item \textbf{CEMP 恒星参数回归数据}：填写观测序列类型（如光谱/测光/其它序列表示）、序列长度与分辨率、目标参数集合（\(P\) 维，如 \(T_{\mathrm{eff}}\)、\(\log g\)、\([\mathrm{Fe/H}]\)、\([\mathrm{C/Fe}]\) 等）、数据清洗与标注来源、以及训练/验证/测试划分原则（如按天区/观测批次/信噪比分层等）。
\end{itemize}

\begin{table}[t]
  \centering
  \caption{数据集统计信息（占位）}
  \label{tab:dataset_stats}
  \begin{tabular}{lcccc}
    \toprule
    数据集/任务 & 频率/分辨率 & 变量数 \(C\) / 参数维 \(P\) & 训练/验证/测试 & 时间跨度/观测范围 \\
    \midrule
    Forecast-Dataset-A & -- & \(C=\) -- & --/--/-- & -- \\
    Forecast-Dataset-B & -- & \(C=\) -- & --/--/-- & -- \\
    CEMP-Regression & -- & \(P=\) -- & --/--/-- & -- \\
    \bottomrule
  \end{tabular}
\end{table}

\subsection{对比方法与骨干设置}
列出基线方法（统计模型、深度模型、Transformer 类等），并说明是否采用相同输入长度、预测长度、特征工程与训练预算。建议给出统一的复现实验设置或引用公开代码设置。

本文渐进式频率条件 FiLM 框架以 iTransformer~\cite{itransformer} 为骨干，默认配置如下（可通过 configs 覆盖）：
\begin{itemize}
  \item \(d_{\text{model}} = 512\)，\(N = 3\) 层 Encoder，\(n_{\text{heads}} = 8\)，\(d_{\text{ff}} = 2048\)；
  \item 频谱参数：\(K = 4\) 频带，3 窗口 \(\mathcal{W} = \{L, L/2, L/4\}\)，latent 维度 \(d_z = 32\)，FiLM 隐层 \(h_f = 32\)；
  \item 调制超参：\(\gamma_{\text{init}} = 0.02\)，\(\delta = 0.7\)，\(s_{\max} = 5.0\)；
  \item 稳定化：\(\log(1+x)\) 压缩、clamp 截断、\(\tanh\) 有界调制、FiLM 末层零初始化——均为默认开启。
\end{itemize}

\subsection{实现细节}
填写优化器（如 AdamW）、学习率、batch size、训练轮数、早停策略、硬件环境与随机种子等。

\subsection{评价指标}
针对两类任务分别报告指标：
\begin{itemize}
  \item \textbf{预测任务}：MAE、MSE/RMSE，并说明是否对变量维与预测步做平均（或报告分 horizon 指标）。
  \item \textbf{参数回归任务}：对每个参数维分别报告 MAE/RMSE（必要时给出加权平均），并可补充 \(R^2\)、Spearman/Pearson 相关系数、以及在不同信噪比区间的分组结果。
\end{itemize}

\subsection{主结果}
\begin{table}[t]
  \centering
  \caption{与对比方法的多步预测结果（占位；越小越好）}
  \label{tab:main_results}
  \begin{tabular}{lcccccccc}
    \toprule
    \multirow{2}{*}{方法} & \multicolumn{2}{c}{\(H=96\)} & \multicolumn{2}{c}{\(H=192\)} & \multicolumn{2}{c}{\(H=336\)} & \multicolumn{2}{c}{\(H=720\)} \\
    \cmidrule(lr){2-3}\cmidrule(lr){4-5}\cmidrule(lr){6-7}\cmidrule(lr){8-9}
    & MSE & MAE & MSE & MAE & MSE & MAE & MSE & MAE \\
    \midrule
    iTransformer (backbone) & -- & -- & -- & -- & -- & -- & -- & -- \\
    Informer~\cite{informer} & -- & -- & -- & -- & -- & -- & -- & -- \\
    Autoformer~\cite{autoformer} & -- & -- & -- & -- & -- & -- & -- & -- \\
    FEDformer~\cite{fedformer} & -- & -- & -- & -- & -- & -- & -- & -- \\
    Progressive Freq.-cond. FiLM (Ours) & \textbf{--} & \textbf{--} & \textbf{--} & \textbf{--} & \textbf{--} & \textbf{--} & \textbf{--} & \textbf{--} \\
    \bottomrule
  \end{tabular}
\end{table}

在文字分析中建议强调：不同预测长度下的趋势、在不同数据集上的一致性、以及与频率结构/非平稳程度相关的表现差异；特别关注相比 iTransformer 骨干的提升幅度，验证频域调制在纯骨干之上的增益。

\begin{table}[t]
  \centering
  \caption{CEMP 恒星参数回归结果（占位；越小越好）}
  \label{tab:cemp_results}
  \begin{tabular}{lcccc}
    \toprule
    方法 & \(T_{\mathrm{eff}}\) MAE & \(\log g\) MAE & \([\mathrm{Fe/H}]\) MAE & \([\mathrm{C/Fe}]\) MAE \\
    \midrule
    StarNet~\cite{starnet} & -- & -- & -- & -- \\
    iTransformer (backbone) & -- & -- & -- & -- \\
    Progressive Freq.-cond. FiLM (Ours) & \textbf{--} & \textbf{--} & \textbf{--} & \textbf{--} \\
    \bottomrule
  \end{tabular}
\end{table}

\subsection{消融实验}
以下消融项逐模块验证渐进式频率条件 FiLM 各核心设计的贡献（按从完整模型逐步移除的顺序）：

\begin{itemize}
  \item \textbf{Full Model}：渐进式频率条件 FiLM 完整模型；
  \item \textbf{$-$ Progressive FiLM}：将逐层独立 FiLM Generator 替换为单一共享 FiLM Generator（仅在第 0 层前调制一次），移除逐层衰减机制，验证"逐层渐进式注入"的价值；
  \item \textbf{$-$ Learned Gamma}：将可学习 \(\gamma_{\text{base}}\) 替换为固定值 0.1（Sigmoid 参数化变为硬编码常数），验证"数据集自适应调制强度"的贡献；
  \item \textbf{$-$ Multi-Resolution}：将三窗口频谱替换为单窗口（仅全序列 \(L\)），验证多分辨率设计的尺度鲁棒性增益；
  \item \textbf{$-$ Spectral Entropy}：从频谱特征中移除归一化频谱熵，仅保留频带能量与高低频比，验证频谱熵作为可预测性信号的价值；
  \item \textbf{$-$ Stabilization}：同时移除 \(\log(1+x)\) 压缩、clamp 截断与 \(\tanh\) 有界约束（调制直接使用原始 \(\alpha, \beta\)），验证稳定化管道对训练鲁棒性的影响；
  \item \textbf{$-$ All Spectrum (= iTransformer)}：完全移除频谱分支与 FiLM 调制，退化为纯 iTransformer + RevIN 基线。
\end{itemize}

\begin{table}[t]
  \centering
  \caption{消融实验结果（占位）}
  \label{tab:ablation}
  \begin{tabular}{lcc}
    \toprule
    变体 & MSE & MAE \\
    \midrule
    Progressive Freq.-cond. FiLM (Full) & \textbf{--} & \textbf{--} \\
    $-$ Progressive FiLM (single FiLM) & -- & -- \\
    $-$ Learned Gamma (fixed \(\gamma = 0.1\)) & -- & -- \\
    $-$ Multi-Resolution (single window) & -- & -- \\
    $-$ Spectral Entropy (energy + ratio only) & -- & -- \\
    $-$ Stabilization (no log1p/clamp/tanh) & -- & -- \\
    $-$ All Spectrum (= iTransformer) & -- & -- \\
    \bottomrule
  \end{tabular}
\end{table}

\subsection{参数敏感性分析（可选）}
建议分析以下关键超参的敏感性：
\begin{itemize}
  \item 衰减因子 \(\delta \in \{0.5, 0.7, 0.85, 1.0\}\)——\(\delta = 1.0\) 等价于所有层等强度调制，可验证衰减机制的必要性；
  \item latent 维度 \(d_z \in \{16, 32, 48, 64\}\)——验证频谱压缩的维度鲁棒性；
  \item 频带数 \(K \in \{3, 4, 6, 8\}\)——验证频带划分粒度的影响；
  \item \(\gamma_{\text{init}}\) 初始值（保守 vs 激进）——验证可退化性初始化策略的有效性。
\end{itemize}

\subsection{可视化分析（可选）}
可包含：(1) 预测曲线与真实曲线的对比；(2) 不同样本的频谱熵与预测误差的关系散点图（验证频谱熵作为可预测性信号的合理性）；(3) 各层 \(\gamma_{\ell}\) 的实际学习值与手选值的对比；(4) 频带能量分布在不同类别序列上的可视化差异。

\begin{figure}[t]
  \centering
  % \includegraphics[width=0.9\linewidth]{figures/placeholder.pdf}
  \fbox{\parbox{0.9\linewidth}{\centering\vspace{3cm}图：预测可视化示例（待补充）}}
  \caption{预测可视化示例（占位）}
  \label{fig:vis}
\end{figure}

% 05 Conclusion
\section{结论}
\label{sec:conclusion}

本文提出\textbf{渐进式频率条件 FiLM}（Progressive Frequency-conditioned FiLM），一种以 iTransformer 为骨干、通过逐层渐进衰减的 FiLM 将多分辨率频谱信息注入 Transformer Encoder 的轻量高效框架。该方法在 iTransformer 的变量维 token 化范式基础上引入三个核心创新：(1) \textbf{多分辨率频谱提取}——在三个时间窗口上通过 FFT 功率谱提取频带归一化能量、高低频比值与频谱熵，形成尺度鲁棒且信息丰富的频域表征；(2) \textbf{逐层渐进式 FiLM}——压缩后的频谱 latent 通过每层独立的小 MLP 生成调制参数，以指数衰减强度注入每一层 Transformer Encoder，实现由粗到细的层级化频域引导；(3) \textbf{可学习调制强度}——通过 Sigmoid 参数化使模型自适应决定频域信息的注入程度，初始保守设置保证可退化性。

在稳定化设计方面，FiLM 末层零初始化、\(\tanh\) 有界约束与 \(\log(1+x)\) 压缩/截断管道共同确保了训练稳健性与可退化性：初始状态下模型等价于 iTransformer + RevIN，频域调制在训练中逐步激活而非骤然介入。框架通过更换任务头可同时覆盖多变量预测（sequence-to-sequence）与参数回归（sequence-to-vector）两类任务。

实验结果与消融分析（待补充）将进一步验证：(1) 多分辨率频谱相比单窗口频谱的增益；(2) 逐层渐进式调制相比单点调制的优势；(3) 可学习 gamma 相比固定系数的数据集自适应性；(4) 各稳定化组件对训练鲁棒性的贡献。

\subsection{局限性与未来工作}
未来可从以下方向进一步拓展：
\begin{itemize}
  \item 探索可学习的窗口尺寸选择或基于注意力/门控的自适应频谱加权，替代固定三窗口设计；
  \item 引入相位信息或复数频谱特征（利用 FFT 的完整信息），进一步丰富频域表征的完整性；
  \item 在多任务设置下联合预测与异常检测/不确定性估计，以增强工业场景可用性；
  \item 探索频谱 latent \(Z\) 的可视化与解释性，建立频谱熵等指标与预测置信度/难度的量化关联；
  \item 将逐层衰减机制从固定指数衰减推广为可学习的衰减曲线（如轻量 MLP 预测每层衰减系数），进一步降低手工超参依赖；
  \item 研究更严格的理论分析（如频域调制对注意力矩阵谱特性的影响），提升方法在分布外（OOD）条件下的可靠性保障。
\end{itemize}


\bibliographystyle{plain}
\bibliography{references}

\appendix
% Appendix (optional)
\section{附录：补充材料（占位）}
\label{sec:appendix}

\subsection{符号表}
\begin{table}[h]
  \centering
  \caption{符号说明}
  \label{tab:symbols}
  \begin{tabular}{ll}
    \toprule
    符号 & 含义 \\
    \midrule
    \(B\) & batch 大小 \\
    \(L\) & 历史输入长度（回望窗口） \\
    \(H\) & 预测步数（预测长度） \\
    \(C\) & 变量数（通道数） \\
    \(P\) & 目标参数维度（回归任务） \\
    \(d_{\text{model}}\) & token 隐空间维度 \\
    \(d_z\) & 频谱 latent 维度 \\
    \(D_s\) & 压缩前频谱特征维度 \\
    \(K\) & 每窗口频带数 \\
    \(\mathcal{W}\) & 多时间窗口集合 \(\{L, L/2, L/4\}\) \\
    \(N\) & Transformer Encoder 层数 \\
    \(\gamma_{\text{base}}\) & 全局可学习调制基强度 \\
    \(\delta\) & 逐层调制衰减因子 \\
    \(\gamma_{\ell}\) & 第 \(\ell\) 层的实际调制强度 \\
    \(\alpha_{\ell}, \beta_{\ell}\) & 第 \(\ell\) 层的 FiLM 缩放/平移参数 \\
    \(M_w\) & 窗口 \(w\) 的功率谱 \\
    \(E_k^{(w)}\) & 窗口 \(w\) 的第 \(k\) 频带能量 \\
    \(R^{(w)}\) & 窗口 \(w\) 的高低频能量比 \\
    \(H^{(w)}\) & 窗口 \(w\) 的归一化频谱熵 \\
    \bottomrule
  \end{tabular}
\end{table}

\subsection{模块计算复杂度分析（可选）}
多分辨率频谱提取的计算复杂度为 \(O(3 \cdot C \cdot L \log L)\)（三次 rFFT，每次在长度为 \(L, L/2, L/4\) 的序列上进行），频谱压缩器与 FiLM Generator 均为轻量 MLP，其计算量显著小于 Transformer Encoder 的自注意力计算 \(O(C^2 \cdot d_{\text{model}})\)。在典型配置下（\(L=96, C=7, d_{\text{model}}=512\)），频谱分支的 FLOPs 仅为 Encoder 部分的约 \(2\%\sim3\%\)。

\subsection{更多结果（可选）}
可在此补充不同数据集、不同 horizon、不同变量维的更完整表格；或补充运行时间/参数量对比等。

\subsection{可退化性验证实验设计（可选）}
建议设计如下验证实验以量化可退化性：
\begin{itemize}
  \item 在纯白噪声（无频谱结构）数据集上，记录 \(\gamma_{\text{base}}\) 在训练过程中的变化曲线——预期其收敛到接近 0 的值；
  \item 在强周期性数据集上，观测 \(\gamma_{\text{base}}\) 是否显著大于初始化值，以及各层 \(\gamma_{\ell}\) 的分布；
  \item 对比"零初始化 vs 随机初始化 FiLM 末层"的训练曲线差异，验证恒等初始化对训练早期稳定性的影响。
\end{itemize}


\end{document}
